\documentclass{article}

\PassOptionsToPackage{numbers,compress}{natbib}

\usepackage{neurips_2026}

\usepackage[utf8]{inputenc}
\usepackage[T1]{fontenc}
\usepackage{hyperref}
\usepackage{url}
\usepackage{amsfonts}
\usepackage{amsmath}
\usepackage{amssymb}
\usepackage{amsthm}
\usepackage{mathtools}
\usepackage{nicefrac}
\usepackage{microtype}
\usepackage{xcolor}
\usepackage{graphicx}

\newcommand{\E}{\mathbb{E}}
\newcommand{\R}{\mathbb{R}}
\newcommand{\diag}{\operatorname{Diag}}

\newtheorem{theorem}{Theorem}
\newtheorem{proposition}[theorem]{Proposition}
\newtheorem{corollary}[theorem]{Corollary}

\title{Designing Diversity in BatchEnsemble}

\author{Anonymous Authors}

\begin{document}

\maketitle

\begin{abstract}
BatchEnsemble is a prominent approach to parameter-efficient ensembling, yet the source and nature of its diversity remain poorly understood. We study BatchEnsemble in a feature-learning infinite-width regime with trainable boundary fast weights and ask a basic structural question: when do ensemble members collapse, and what deterministic baseline does that collapse produce? We show that jointly symmetric initialization leads to exact functional collapse even when the shared representation and the boundary fast weights are trained. The resulting collapsed system is not the ordinary single-model limit: the boundary variables follow the same single-member drift at a fixed state, while the shared hidden-neuron law is accelerated by the ensemble width. This identifies the correct deterministic reference point for analyzing diversity in parameter-efficient ensembles.
\end{abstract}

\section{Introduction}

Deep ensembles remain one of the most reliable tools for improving both predictive performance and predictive uncertainty, but their training and inference costs scale roughly linearly with the number of members \cite{lakshminarayanan2017}. This cost motivates the study of \emph{efficient ensembles}: architectures that preserve ensemble-like behavior while sharing most of the computation and memory. BatchEnsemble is one of the most prominent examples in this direction. Starting from a standard neural network, BatchEnsemble keeps the main weight matrices shared across all members and gives each member only a small number of multiplicative fast weights. Concretely, for a dense layer with shared weight matrix $W$ and ensemble member $\alpha\in\{1,\dots,K\}$, the member-specific matrix is
\[
W_\alpha = W \odot (r_\alpha s_\alpha^\top),
\]
where $r_\alpha$ and $s_\alpha$ are lightweight member-specific vectors and $\odot$ denotes the Hadamard product \cite{wen2020}. Thus, one can think of BatchEnsemble as an ensemble of subnetworks that share the bulk of their parameters but differ through low-rank multiplicative gates applied to each shared layer. This structure makes it possible to train and evaluate several predictors at only modest overhead compared with a single network.

Parameter-efficient ensembling has recently become an important primitive in tabular deep learning through TabM, an MLP-based architecture built around efficient ensembling that reports the best performance among tabular deep learning models on a large public benchmark, remains among the strongest individual deep tabular models in subsequent benchmark studies, and has also appeared in top Kaggle solutions \cite{gorishniy2025tabm,tabarena2025,gorishniy2025tabmrepo}. More broadly, if a cheap ensemble retains meaningful member diversity, then one can hope to obtain better calibrated predictions and richer uncertainty estimates without paying the full price of a deep ensemble. If, on the other hand, its members collapse toward nearly identical predictors, then the promise of parameter-efficient ensembling becomes much less clear.

At the same time, the existing picture is mixed. On the empirical side, some studies report that BatchEnsemble-like methods can provide uncertainty estimates competitive with much more expensive ensembles \cite{blorstad2026}, while others argue that BatchEnsemble members can behave much more like a single model than like a genuine ensemble \cite{zamyatin2026}. On the theoretical side, direct analysis of BatchEnsemble-style efficient ensembles remains limited, with the main asymptotic results concentrating on NTK/lazy regimes \cite{jacot2018,velikanov2022}. Here, \emph{lazy} refers to settings in which the parameters remain close to their initialization throughout training, so the network behaves approximately like a fixed kernel model, whereas \emph{feature learning} refers to the more practical setting in which the parameters can move substantially away from initialization and the learned representation itself changes during training \cite{yanghu2020}. A corresponding theory for this feature-learning setting with trainable boundary fast weights is still missing.

In this paper, we study this question through a mean-field analysis centered on the symmetric reference point. We show that symmetric initialization leads to collapse even when the shared representation and the boundary fast weights are trainable. However, the resulting collapsed predictor is not an ordinary single-model limit: it follows a distinct shared-law dynamics that depends on the ensemble width. This isolates the correct deterministic baseline against which any symmetry-breaking mechanism must be compared.

\paragraph{Contributions.}
Our contributions are summarized as follows:
\begin{itemize}
\item We establish a symmetric-collapse theory for BatchEnsemble in a feature-learning infinite-width regime with trainable boundary fast weights, and identify the corresponding collapsed baseline.
\item We prove that the collapsed baseline is not the matched ordinary single-model flow: only the shared hidden-neuron law is multiplied by a factor $K$, while the boundary variables keep the same single-member drift at a fixed state.
\end{itemize}

% TODO: Add citations in the final draft.
% If space becomes tight, the repository citation can be moved to a footnote.

\section{Preliminary}

\subsection{Formulation}

We study a two-layer BatchEnsemble-style predictor
\begin{equation}
f^M_\alpha(x)
:=
\frac{1}{M}
\sum_{m=1}^M
(a_m\odot r_\alpha^{(2)})\,s^{(2)}_{\alpha,m}\,
\sigma\!\bigl(r^{(1)}_{\alpha,m}\, w_m^\top(s^{(1)}_\alpha\odot x)\bigr),
\label{eq:full-predictor}
\end{equation}
where $x\in\R^{d_x}$ is the input, $f^M_\alpha(x)\in\R^{d_y}$ is the prediction of ensemble member $\alpha\in\{1,\dots,K\}$, and $m\in\{1,\dots,M\}$ indexes hidden neurons. The vectors $a_m\in\R^{d_y}$ and $w_m\in\R^{d_x}$ are shared output- and input-side neuron parameters, respectively. The scalars $s^{(2)}_{\alpha,m}$ and $r^{(1)}_{\alpha,m}$ are neuron-wise fast weights, while $s^{(1)}_\alpha\in\R^{d_x}$ and $r^{(2)}_\alpha\in\R^{d_y}$ are boundary fast weights acting directly on the input and output sides. Finally, $\sigma$ is the activation function and $\odot$ denotes elementwise multiplication. This decomposition separates two qualitatively different sources of member-specific variation: neuron-wise fast weights $\{s^{(2)}_{\alpha,m},r^{(1)}_{\alpha,m}\}$, which are attached to individual hidden units, and boundary fast weights $\{s^{(1)}_\alpha,r^{(2)}_\alpha\}$, which act directly on the input or output side of the predictor.

The main quantities of interest are the mean predictor
\[
\bar f^M(x):=\frac{1}{K}\sum_{\alpha=1}^K f^M_\alpha(x),
\]
the disagreement or predictive spread
\[
\mathrm{Dis}^M(x):=\frac{1}{K}\sum_{\alpha=1}^K \|f^M_\alpha(x)-\bar f^M(x)\|^2,
\]
and the ensemble risk induced by $\bar f^M$,
\[
\mathcal R(\bar f^M):=\E_{(x,y)\sim \mathcal D}\bigl[\ell(\bar f^M(x),y)\bigr].
\]

\subsection{Related Work}

This paper sits at the intersection of efficient ensembling, infinite-width theory, and uncertainty-oriented analysis of BatchEnsemble. On the modeling side, BatchEnsemble was introduced by Wen et al.\ as a parameter-sharing alternative to deep ensembles, with the goal of preserving accuracy and predictive uncertainty benefits at much lower cost \cite{wen2020}. Our work does not propose a new architecture; rather, it asks what kinds of member diversity survive when such an architecture is taken to a feature-learning infinite-width limit.

On the theory side, it is useful to distinguish two standard infinite-width viewpoints. In the neural tangent kernel (NTK) or lazy regime, one linearizes the network around its initialization and obtains a deterministic kernel dynamics governed by the NTK \cite{jacot2018}. This regime is important because it gives a clean asymptotic theory for shared-parameter models without requiring substantial movement away from initialization. In contrast, mean-field or feature-learning analyses track an evolving population law for the hidden units, so the limiting predictor depends on the learned population state rather than only on the initialization kernel. Following Yang and Hu, we view these as genuinely different asymptotic regimes rather than two equivalent descriptions of the same wide-network limit \cite{yanghu2020}.

Among works closest to BatchEnsemble-style models, Velikanov et al.\ analyze Embedded Ensembles in the NTK regime and show that wide shared-parameter ensembles can display either \emph{collective} behavior, in which members follow a common kernel dynamics, or \emph{independent} behavior, in which different members retain distinct limiting predictors \cite{velikanov2022}. This result is especially relevant for us because it shows that even within infinite-width shared-parameter models, qualitatively different kinds of diversity are possible. However, the analysis remains in a lazy-training regime. Our contribution is to identify the corresponding picture in a feature-learning setting with trainable boundary fast weights, where the central issues become collapse, surviving boundary asymmetry, and the design of useful diversity.

Finally, our work is motivated by mixed empirical evidence on BatchEnsemble diversity. Bl\o rstad et al.\ report strong uncertainty-estimation performance for BatchEnsemble across tabular and time-series tasks \cite{blorstad2026}, whereas Zamyatin et al.\ argue that BatchEnsemble can behave much more like a single model than like a genuine ensemble \cite{zamyatin2026}. We interpret this tension as evidence that the relevant question is not whether BatchEnsemble is ``diverse'' in the abstract, but which asymmetries survive at leading order and whether they move the central predictor, increase predictive spread, or both.

\section{Theoretical Results}

Our theoretical analysis concerns the infinite-width limits of the formulation above. Taking $M\to\infty$ yields a feature-learning mean-field description in terms of a hidden-neuron law and, when boundary fast weights are present, a finite-dimensional member state built from those fast weights. In this section we focus on the symmetric reference point and identify the deterministic baseline that it induces.

For the collapse theory, it is convenient to introduce a \emph{member-state predictor}. If
\[
u\in\R^{d_x},
\qquad
v\in\R^{d_y}
\]
denote the input- and output-side boundary fast weights for one member, we write
\begin{equation}
H(x;\rho,u,v)
:=
\int_{\R^{d_y}\times\R^{d_x}}
(\tilde a\odot v)\,
\sigma\!\bigl(\tilde w^\top(u\odot x)\bigr)\,
\rho(\mathrm{d}\tilde a,\mathrm{d}\tilde w),
\label{eq:member-state-predictor}
\end{equation}
Here $\rho$ is the common infinite-width neuron law. The integration variables $(\tilde a,\tilde w)$ under $\rho$ are not the finite-width shared parameters $(a_m,w_m)$ from \eqref{eq:full-predictor}; they should be read as effective hidden-neuron coordinates after any neuron-wise member-specific randomness has self-averaged into the common law. The point of $H$ is that it keeps the input- and output-side member-specific coordinates $(u,v)$ explicit while packaging the remaining hidden-neuron randomness into $\rho$. It is not a separate finite-width architecture, and it does not by itself specify whether these coordinates are initialized symmetrically, randomized, trainable, or frozen. Those regime assumptions are stated directly in each subsection in terms of the parameters of \eqref{eq:full-predictor}.

\subsection{Symmetric Initialization and the Collapsed Baseline}

Our first result identifies the symmetric reference point. In the feature-learning regime considered here, jointly permutation-symmetric initialization still leads to collapse, even though the shared representation, the neuron-wise fast weights, and the boundary fast weights are all trainable.

\paragraph{Setup for this subsection.}
We work with the predictor \eqref{eq:full-predictor}. All parameter families may be trained: the shared weights $\{a_m,w_m\}_{m=1}^M$, the neuron-wise fast weights $\{s^{(2)}_{\alpha,m},r^{(1)}_{\alpha,m}\}_{\alpha,m}$, and the boundary fast weights $\{s^{(1)}_\alpha,r^{(2)}_\alpha\}_{\alpha=1}^K$. The initialization assumption is only that the boundary coordinates are identical across members at time $0$ and that the empirical law of the neuron-level tuples
\[
\bigl(a_m,w_m,\{s^{(2)}_{\alpha,m},r^{(1)}_{\alpha,m}\}_{\alpha=1}^K\bigr)
\]
converges to a member-permutation-symmetric limit.

To state the proof cleanly, write
\[
q_\alpha:=(u_\alpha,v_\alpha)\in\R^{d_x}\times\R^{d_y},
\qquad
q:=(q_1,\dots,q_K).
\]
We assume the width-limit dynamics close into a deterministic well-posed mean-field system
\[
\dot\rho_t=\mathfrak B(\rho_t,q_t),
\qquad
\dot q_{\alpha,t}=\mathfrak C_\alpha(\rho_t,q_t),
\qquad
\alpha=1,\dots,K,
\]
where $\rho_t$ is a law on hidden-neuron coordinates that may still carry the head-labeled neuron-wise fast weights. For each permutation $\pi\in S_K$, let $\Pi_\pi\rho$ denote the pushforward of $\rho$ by permuting those head labels inside one neuron state, and let
\[
\Pi_\pi q:=(q_{\pi^{-1}(1)},\dots,q_{\pi^{-1}(K)}).
\]
Because the population objective is the head-summed loss induced by \eqref{eq:full-predictor}, the predictor map and the limiting vector field are permutation-equivariant:
\begin{align}
f_\alpha(x;\Pi_\pi\rho,\Pi_\pi q)
&=
f_{\pi^{-1}(\alpha)}(x;\rho,q),
\label{eq:paper-perm-pred}\\
\mathfrak B(\Pi_\pi\rho,\Pi_\pi q)
&=
\Pi_\pi\mathfrak B(\rho,q),
\label{eq:paper-perm-rho}\\
\mathfrak C_\alpha(\Pi_\pi\rho,\Pi_\pi q)
&=
\mathfrak C_{\pi^{-1}(\alpha)}(\rho,q).
\label{eq:paper-perm-q}
\end{align}
These are the only structural facts needed below.

\begin{theorem}[Symmetric initialization yields wide-limit collapse]
\label{thm:paper-collapse}
Let $\rho_t$ denote the deterministic hidden-neuron law and let
\[
(u_{\alpha,t},v_{\alpha,t})\in\R^{d_x}\times\R^{d_y},
\qquad
\alpha=1,\dots,K,
\]
denote the corresponding member-specific input- and output-side boundary states. If the initial hidden-neuron law and these member states are jointly permutation-symmetric across the $K$ members, then this symmetry is preserved for all $t\ge 0$, and consequently
\[
f_1(x;\rho_t,u_{1,t},v_{1,t})=\cdots=f_K(x;\rho_t,u_{K,t},v_{K,t})
\qquad
\forall x,\ \forall t\ge 0.
\]
\end{theorem}

\begin{proof}
Fix a permutation $\pi\in S_K$ and define the transformed trajectory
\[
\widetilde\rho_t:=\Pi_\pi\rho_t,
\qquad
\widetilde q_t:=\Pi_\pi q_t.
\]
If $(\rho_t,q_t)$ solves the limiting system, then \eqref{eq:paper-perm-rho} and \eqref{eq:paper-perm-q} imply that $(\widetilde\rho_t,\widetilde q_t)$ solves the same system with initial condition
\[
(\widetilde\rho_0,\widetilde q_0)
=
(\Pi_\pi\rho_0,\Pi_\pi q_0).
\]
By the joint permutation symmetry of the initial hidden-neuron law and the initial member states,
\[
(\Pi_\pi\rho_0,\Pi_\pi q_0)
=
(\rho_0,q_0).
\]
Uniqueness of the deterministic mean-field evolution therefore yields
\[
\Pi_\pi\rho_t=\rho_t,
\qquad
\Pi_\pi q_t=q_t
\qquad
\forall t\ge 0.
\]
Since this holds for every permutation, all member states coincide along the trajectory:
\[
q_{1,t}=\cdots=q_{K,t}
\qquad
\forall t\ge 0.
\]
Now fix any pair of indices $\alpha,\beta\in\{1,\dots,K\}$ and choose $\pi$ so that $\pi^{-1}(\alpha)=\beta$. Using \eqref{eq:paper-perm-pred} together with the invariance just proved,
\[
f_\alpha(x;\rho_t,q_t)
=
f_\alpha(x;\Pi_\pi\rho_t,\Pi_\pi q_t)
=
f_{\pi^{-1}(\alpha)}(x;\rho_t,q_t)
=
f_\beta(x;\rho_t,q_t).
\]
Because $\alpha$ and $\beta$ were arbitrary, all member predictors agree for every $x$ and every $t\ge 0$.
\end{proof}

In the concrete BatchEnsemble model, this means that all members start with the same input-side and output-side boundary vectors, while the full initial neuron law---including any neuron-wise fast-weight randomness---converges to a permutation-symmetric limit across members. Random iid neuron-wise initializations are therefore covered as long as their large-width empirical law is head-exchangeable, and the no-boundary architecture is the special case in which these common boundary vectors are fixed to the all-ones vectors.

This statement is already nontrivial. It says that collapse persists in a feature-learning setting where the backbone, the neuron-wise fast weights, and the boundary variables are all trained; it is therefore not a lazy-training artifact of a frozen representation. The practical message is that simply making boundary fast weights trainable does not create order-one diversity if the infinite-width state is initialized on the permutation-symmetric diagonal.

The second point is equally important: the common predictor obtained in Theorem~\ref{thm:paper-collapse} is not an ordinary single-model limit. The proposition below identifies the dynamics of that common predictor.

Under the assumptions of Theorem~\ref{thm:paper-collapse}, write
\[
u_t:=u_{1,t}=\cdots=u_{K,t},
\qquad
v_t:=v_{1,t}=\cdots=v_{K,t},
\]
and let $\mathcal B(\rho,u,v)$ denote the shared-law drift of the matched ordinary single-model mean-field dynamics, while $\mathcal C_u(\rho,u,v)$ and $\mathcal C_v(\rho,u,v)$ denote the drifts of the input- and output-side boundary variables. At the level of the limit system, these are defined by the head-additive/head-local decomposition
\begin{equation}
\mathfrak B(\rho,q_1,\dots,q_K)
=
\sum_{\alpha=1}^K \mathcal B(\rho,q_\alpha),
\qquad
\mathfrak C_\alpha(\rho,q_1,\dots,q_K)
=
\bigl(\mathcal C_u(\rho,q_\alpha),\mathcal C_v(\rho,q_\alpha)\bigr).
\label{eq:paper-additive-local}
\end{equation}
This is exactly the structure induced by the head-summed loss: the shared hidden variables see a sum over heads, while differentiating with respect to $(u_\alpha,v_\alpha)$ removes all terms except head $\alpha$.

\begin{proposition}[The collapsed baseline is a $K$-accelerated shared-law flow]
\label{prop:paper-collapse-baseline}
The common collapsed predictor from Theorem~\ref{thm:paper-collapse} takes the form
\begin{align*}
f_{\mathrm{coll}}(x,t)&=H(x;\rho_t,u_t,v_t),\\
(\dot\rho_t,\dot u_t,\dot v_t)&=\bigl(K\,\mathcal B(\rho_t,u_t,v_t),\ \mathcal C_u(\rho_t,u_t,v_t),\ \mathcal C_v(\rho_t,u_t,v_t)\bigr),
\end{align*}
whereas the matched ordinary single-model baseline uses the vector field
\[
(\dot\rho_t,\dot u_t,\dot v_t)=\bigl(\mathcal B(\rho_t,u_t,v_t),\ \mathcal C_u(\rho_t,u_t,v_t),\ \mathcal C_v(\rho_t,u_t,v_t)\bigr).
\]
Thus, at the same state $(\rho,u,v)$, the two systems have exactly the same drift in the input- and output-side boundary variables and differ only by a factor $K$ in the drift of the shared hidden-neuron law. Their trajectories therefore generally differ, but the difference is entirely on the shared-law side.
\end{proposition}

\begin{proof}
By Theorem~\ref{thm:paper-collapse}, there exists a common boundary path
\[
q_t=(u_t,v_t)
\]
such that
\[
q_{1,t}=\cdots=q_{K,t}=q_t
\qquad
\forall t\ge 0.
\]
Substituting this diagonal trajectory into \eqref{eq:paper-additive-local} gives
\[
\dot\rho_t
=
\sum_{\alpha=1}^K \mathcal B(\rho_t,q_t)
=
K\,\mathcal B(\rho_t,q_t),
\]
and
\[
\dot q_t
=
\bigl(\mathcal C_u(\rho_t,q_t),\mathcal C_v(\rho_t,q_t)\bigr).
\]
Writing $q_t=(u_t,v_t)$ yields
\[
(\dot\rho_t,\dot u_t,\dot v_t)
=
\bigl(
K\,\mathcal B(\rho_t,u_t,v_t),\,
\mathcal C_u(\rho_t,u_t,v_t),\,
\mathcal C_v(\rho_t,u_t,v_t)
\bigr).
\]
The common predictor is therefore
\[
f_{\mathrm{coll}}(x,t)=H(x;\rho_t,u_t,v_t).
\]
By definition, the matched ordinary single-model baseline evolves with a single copy of the shared-law drift and the same boundary drifts at the same state, namely
\[
(\dot\rho_t,\dot u_t,\dot v_t)
=
\bigl(
\mathcal B(\rho_t,u_t,v_t),\,
\mathcal C_u(\rho_t,u_t,v_t),\,
\mathcal C_v(\rho_t,u_t,v_t)
\bigr).
\]
Comparing the two systems proves the claim.
\end{proof}

So the correct deterministic benchmark for randomized BatchEnsemble is the \emph{collapsed BatchEnsemble baseline}, not an ordinary single model.

\subsection{Linearized Symmetry Breaking: Consensus and Contrast Modes}

Proposition~\ref{prop:paper-collapse-baseline} identifies the correct deterministic reference trajectory on the permutation-symmetric diagonal. The next question is how small asymmetries evolve around that diagonal, and which ones generate useful disagreement rather than merely perturbing the common predictor.

Write the collapsed trajectory as
\[
(\rho_t^\star,q_t^\star)
\qquad
\text{with}
\qquad
q_t^\star=(u_t^\star,v_t^\star),
\]
so that
\[
\dot\rho_t^\star = K\,\mathcal B(\rho_t^\star,q_t^\star),
\qquad
\dot q_t^\star = \mathcal C(\rho_t^\star,q_t^\star),
\]
where $\mathcal C:=(\mathcal C_u,\mathcal C_v)$. Assume for this subsection that $\mathcal B$, $\mathcal C$, and the predictor map $H$ from \eqref{eq:member-state-predictor} are continuously differentiable along this path. Denote the corresponding derivatives by
\[
B_\rho(t):=D_\rho\mathcal B(\rho_t^\star,q_t^\star),
\quad
B_q(t):=D_q\mathcal B(\rho_t^\star,q_t^\star),
\]
\[
C_\rho(t):=D_\rho\mathcal C(\rho_t^\star,q_t^\star),
\quad
C_q(t):=D_q\mathcal C(\rho_t^\star,q_t^\star),
\]
and
\[
J_\rho(t,x):=D_\rho H(x;\rho_t^\star,q_t^\star),
\qquad
J_q(t,x):=D_q H(x;\rho_t^\star,q_t^\star).
\]

Consider a perturbation of the full $K$-member system of the form
\[
\rho_t = \rho_t^\star + \varepsilon \eta_t + o(\varepsilon),
\qquad
q_{\alpha,t} = q_t^\star + \varepsilon \zeta_{\alpha,t} + o(\varepsilon).
\]
Decompose the member perturbations into their empirical mean and centered part:
\[
\bar\zeta_t:=\frac{1}{K}\sum_{\alpha=1}^K \zeta_{\alpha,t},
\qquad
\xi_{\alpha,t}:=\zeta_{\alpha,t}-\bar\zeta_t,
\qquad
\sum_{\alpha=1}^K \xi_{\alpha,t}=0.
\]

\begin{theorem}[Consensus/contrast decomposition around the collapsed diagonal]
\label{thm:paper-consensus-contrast}
The linearized dynamics split into two invariant sectors:
\begin{align}
\dot\eta_t
&=
K\,B_\rho(t)\eta_t + K\,B_q(t)\bar\zeta_t,
\label{eq:paper-lin-rho}\\
\dot{\bar\zeta}_t
&=
C_\rho(t)\eta_t + C_q(t)\bar\zeta_t,
\label{eq:paper-lin-mean}\\
\dot\xi_{\alpha,t}
&=
C_q(t)\xi_{\alpha,t},
\qquad
\alpha=1,\dots,K.
\label{eq:paper-lin-contrast}
\end{align}
Moreover, the member predictors satisfy
\begin{align}
f_\alpha(x,t)-\bar f(x,t)
=
\varepsilon J_q(t,x)\xi_{\alpha,t} + o(\varepsilon),
\label{eq:paper-lin-disagreement}
\end{align}
whereas the ensemble mean satisfies
\begin{align}
\bar f(x,t)
=
H(x;\rho_t^\star,q_t^\star)
\,+\,
\varepsilon\bigl[J_\rho(t,x)\eta_t + J_q(t,x)\bar\zeta_t\bigr]
\,+\,
o(\varepsilon).
\label{eq:paper-lin-mean-predictor}
\end{align}
Hence only the centered contrast modes $\{\xi_{\alpha,t}\}_{\alpha=1}^K$ contribute to disagreement at first order, while the consensus mode $(\eta_t,\bar\zeta_t)$ controls the first-order shift of the common predictor.
\end{theorem}

\begin{proof}
Linearizing the full system
\[
\dot\rho_t = \sum_{\alpha=1}^K \mathcal B(\rho_t,q_{\alpha,t}),
\qquad
\dot q_{\alpha,t} = \mathcal C(\rho_t,q_{\alpha,t})
\]
around the diagonal trajectory $(\rho_t^\star,q_t^\star,\dots,q_t^\star)$ gives
\[
\dot\eta_t
=
\sum_{\alpha=1}^K \bigl(B_\rho(t)\eta_t + B_q(t)\zeta_{\alpha,t}\bigr)
=
K\,B_\rho(t)\eta_t + K\,B_q(t)\bar\zeta_t,
\]
and
\[
\dot\zeta_{\alpha,t}
=
C_\rho(t)\eta_t + C_q(t)\zeta_{\alpha,t}.
\]
Averaging the second display over $\alpha$ yields \eqref{eq:paper-lin-mean}, and subtracting the averaged equation from the $\alpha$th equation yields \eqref{eq:paper-lin-contrast}. Linearizing the predictor map gives
\[
f_\alpha(x,t)
=
H(x;\rho_t^\star,q_t^\star)
\,+\,
\varepsilon\bigl[J_\rho(t,x)\eta_t + J_q(t,x)\zeta_{\alpha,t}\bigr]
\,+\,
o(\varepsilon),
\]
from which \eqref{eq:paper-lin-disagreement} and \eqref{eq:paper-lin-mean-predictor} follow after writing $\zeta_{\alpha,t}=\bar\zeta_t+\xi_{\alpha,t}$ and averaging.
\end{proof}

Theorem~\ref{thm:paper-consensus-contrast} has a sharp design consequence. If one wants disagreement without perturbing the common predictor, then the perturbation budget should be placed entirely in the centered contrast sector.

Let $\Phi_C(t,s)$ denote the fundamental solution of \eqref{eq:paper-lin-contrast}, let
\[
\Sigma_0
:=
\frac{1}{K}\sum_{\alpha=1}^K \xi_{\alpha,0}\xi_{\alpha,0}^\top,
\]
and define the propagated contrast Gram matrix
\[
G_t(x)
:=
\Phi_C(t,0)^\top J_q(t,x)^\top J_q(t,x)\Phi_C(t,0).
\]

\begin{corollary}[Leading disagreement depends only on contrast covariance]
\label{cor:paper-contrast-cov}
If the initial perturbation is centered in the sense that
\[
\eta_0=0,
\qquad
\bar\zeta_0=0,
\]
then
\begin{equation}
\mathrm{Dis}(x,t)
=
\varepsilon^2 \operatorname{tr}\!\bigl(G_t(x)\Sigma_0\bigr)
\,+\,
o(\varepsilon^2).
\label{eq:paper-contrast-cov}
\end{equation}
Thus the leading disagreement depends on the initialization only through the empirical covariance of the contrast codes.
\end{corollary}

\begin{proof}
Under $\eta_0=0$ and $\bar\zeta_0=0$, the homogeneous system \eqref{eq:paper-lin-rho}--\eqref{eq:paper-lin-mean} implies $\eta_t\equiv 0$ and $\bar\zeta_t\equiv 0$. Hence
\[
\xi_{\alpha,t}=\Phi_C(t,0)\xi_{\alpha,0}
\]
and
\[
f_\alpha(x,t)-\bar f(x,t)
=
\varepsilon J_q(t,x)\Phi_C(t,0)\xi_{\alpha,0} + o(\varepsilon).
\]
Substituting this into the definition of $\mathrm{Dis}(x,t)$ yields \eqref{eq:paper-contrast-cov}.
\end{proof}

Equation~\eqref{eq:paper-contrast-cov} turns small-budget diversity design into a covariance-design problem. In particular, under the trace constraint
\[
\operatorname{tr}(\Sigma_0)\le \sigma^2,
\]
the optimal leading-order disagreement coefficient is $\sigma^2\lambda_{\max}(G_t(x))$ and is achieved by a rank-one antithetic code: half the heads are initialized with $+\sigma v_{\max}$ and half with $-\sigma v_{\max}$, where $v_{\max}$ is a top eigenvector of $G_t(x)$. This shows that, near the collapsed diagonal, iid isotropic randomness is generally suboptimal, and increasing $K$ by itself does not increase the best achievable first-order disagreement coefficient.

Furthermore, Theorem~\ref{thm:paper-consensus-contrast} implies that centered contrast perturbations do not move the ensemble mean predictor at order $\varepsilon$. Consequently, under the additional mild assumption that the solution map and the loss are $C^2$ in the perturbation scale, the ensemble risk of the mean predictor changes only at order $O(\varepsilon^2)$. So, to leading order, one can increase disagreement without paying a first-order price in the central predictor.


\section{Experiments}

Empirical validation is left for future work. The most direct experiment suggested by Section~3 is to compare the ordinary single-model trajectory with the collapsed BatchEnsemble baseline under symmetric boundary initialization and to test whether the factor $K$ in the shared-law update already creates a visible training-dynamics gap before any explicit symmetry breaking is introduced.

\section{Conclusion}

We studied BatchEnsemble through the lens of symmetric initialization in a feature-learning infinite-width regime. Our analysis shows that jointly permutation-symmetric initialization leads to exact collapse, but also that the resulting collapsed predictor is not merely an ordinary single-model baseline. The correct deterministic reference object is a collapsed BatchEnsemble flow in which the boundary variables evolve with the same single-member drift while the shared hidden-neuron law is accelerated by a factor $K$. We view this as the structural baseline that any future theory of symmetry breaking or useful diversity should build on.

% TODO: In the final draft, close with concrete future directions:
% - order-one input-side randomness
% - stronger training-map sensitivity theory
% - broader architecture classes

\section*{References}

\begin{thebibliography}{99}
\bibitem{lakshminarayanan2017}
Balaji Lakshminarayanan, Alexander Pritzel, and Charles Blundell.
\newblock Simple and scalable predictive uncertainty estimation using deep ensembles.
\newblock In \emph{Advances in Neural Information Processing Systems}, 2017.

\bibitem{wen2020}
Yeming Wen, Dustin Tran, and Jimmy Ba.
\newblock BatchEnsemble: An alternative approach to efficient ensemble and lifelong learning.
\newblock In \emph{International Conference on Learning Representations}, 2020.

\bibitem{gorishniy2025tabm}
Yury Gorishniy, Akim Kotelnikov, and Artem Babenko.
\newblock TabM: Advancing tabular deep learning with parameter-efficient ensembling.
\newblock In \emph{International Conference on Learning Representations}, 2025.

\bibitem{gorishniy2025tabmrepo}
Yandex Research.
\newblock TabM: Advancing tabular deep learning with parameter-efficient ensembling.
\newblock GitHub repository, \url{https://github.com/yandex-research/tabm}, accessed March 29, 2026.

\bibitem{tabarena2025}
Nick Erickson, Lennart Purucker, Andrej Tschalzev, David Holzm\"uller, Prateek Mutalik Desai, David Salinas, and Frank Hutter.
\newblock TabArena: A living benchmark for machine learning on tabular data.
\newblock arXiv preprint arXiv:2506.16791, 2025.

\bibitem{blorstad2026}
Morten Bl{\o}rstad, Herman Jangsett Mostein, Nello Blaser, and Pekka Parviainen.
\newblock Evaluating prediction uncertainty estimates from BatchEnsemble.
\newblock arXiv:2601.21581, 2026.

\bibitem{zamyatin2026}
Anton Zamyatin, Patrick Indri, Sagar Malhotra, and Thomas G{\"a}rtner.
\newblock Is BatchEnsemble a single model? On calibration and diversity of efficient ensembles.
\newblock arXiv:2601.16936, 2026.

\bibitem{jacot2018}
Arthur Jacot, Franck Gabriel, and Cl{\'e}ment Hongler.
\newblock Neural tangent kernel: Convergence and generalization in neural networks.
\newblock In \emph{Advances in Neural Information Processing Systems}, 2018.

\bibitem{yanghu2020}
Greg Yang and Edward J. Hu.
\newblock Feature learning in infinite-width neural networks.
\newblock arXiv:2011.14522, 2020.

\bibitem{velikanov2022}
Maksim Velikanov, Roman V. Kail, Ivan Anokhin, Roman Vashurin, Maxim Panov, Alexey Zaytsev, and Dmitry Yarotsky.
\newblock Embedded ensembles: Infinite width limit and operating regimes.
\newblock In \emph{Proceedings of The 25th International Conference on Artificial Intelligence and Statistics}, pages 3138--3163, 2022.
\end{thebibliography}

\end{document}
